% 图 1.2 GAN-based JSCC 结构图 — 基于 CICC0903540 技术文档 2.2
\documentclass[tikz,border=8pt]{standalone}
% 顶会/顶刊风格：低饱和度马卡龙 + 高级感
% 适用于 NeurIPS, CVPR, ICLR, IEEE TPAMI
\usepackage{amssymb}
\usepackage{fontspec}
\usepackage{xeCJK}
\setCJKmainfont{SimHei}
\usepackage{tikz}
\usepackage{pgfplots}
\usetikzlibrary{
  arrows.meta, positioning, shapes.geometric, calc,
  backgrounds, fit, decorations.pathreplacing, patterns
}

% 低饱和度马卡龙配色
\definecolor{mBlue}{HTML}{AEC6CF}
\definecolor{mPink}{HTML}{FFD1DC}
\definecolor{mGreen}{HTML}{B1D8B7}
\definecolor{mPurple}{HTML}{B39EB5}
\definecolor{mYellow}{HTML}{FDFD96}
\definecolor{mGray}{HTML}{E5E4E2}
\definecolor{mDark}{HTML}{4A4A4A}

% 全局 TikZ 样式
\tikzset{
  >=stealth,
  every node/.style={align=center, font=\sffamily},
  block/.style={
    rounded corners=3pt, draw=mDark, align=center,
    minimum height=0.8cm, inner sep=5pt},
  arr/.style={-{Stealth[length=4pt]}, semithick, draw=mDark},
  % 信息密度增强：张量维度标注
  ts/.style={font=\tiny\sffamily, color=mDark!80},
  % 数学操作符节点（⊕加 ⊗乘 ‖拼接）
  op/.style={circle, draw=mDark, minimum size=0.4cm, inner sep=0,
    font=\scriptsize\sffamily, fill=mGray!60},
  % 图例
  legendbox/.style={rectangle, draw=mDark!50, rounded corners=2pt,
    fill=mGray!15, inner sep=4pt, font=\tiny\sffamily},
}

\begin{document}
\begin{tikzpicture}[
  node distance=0.6cm and 0.8cm,
  block/.style={
    rectangle, rounded corners=3pt, draw=mDark, align=center,
    font=\small\sffamily, semithick,
    minimum width=1.7cm, minimum height=0.75cm, inner sep=4pt},
  io/.style={
    rectangle, rounded corners=3pt, draw=mDark, align=center,
    font=\small\sffamily, semithick,
    minimum width=1.2cm, minimum height=0.65cm, inner sep=3pt},
  ch/.style={
    rectangle, rounded corners=3pt, draw=mDark, align=center,
    font=\small\sffamily, semithick,
    minimum width=1.5cm, minimum height=0.75cm, inner sep=3pt},
  disc/.style={
    rectangle, rounded corners=3pt, draw=mDark, align=center,
    font=\small\sffamily, semithick, dashed,
    minimum width=1.8cm, minimum height=0.7cm, inner sep=3pt},
  micro/.style={
    rectangle, rounded corners=2pt, draw=mDark!70, align=center,
    font=\tiny\sffamily, inner sep=2pt, minimum height=0.35cm},
  arr/.style={-{Stealth[length=3pt]}, draw=mDark, semithick},
  darr/.style={-{Stealth[length=3pt]}, draw=mDark!70, semithick, dashed},
]
  % === 主流程（放宽横向间距）===
  \node[io, fill=mGray!50] (img) {$\mathbf{x}$};
  \node[block, fill=mBlue!40, right=0.8cm of img] (enc) {Encoder};
  \node[block, fill=mBlue!40, right=0.8cm of enc] (pwr) {功率控制};
  \node[ch, fill=mGreen!50, right=0.8cm of pwr] (awgn) {AWGN};
  \node[block, fill=mPurple!40, right=0.8cm of awgn] (gen) {Generator};
  \node[io, fill=mGray!50, right=0.8cm of gen] (out) {$\hat{\mathbf{x}}$};

  % 张量标注（放在箭头上方，避免与下方重叠）
  \draw[arr] (img) -- (enc)
    node[pos=0.5, above=2pt, ts] {$1{\times}3{\times}256{\times}256$};
  \draw[arr] (enc) -- (pwr)
    node[pos=0.5, above=2pt, ts] {$\mathbf{y}$};
  \draw[arr] (pwr) -- (awgn);
  \draw[arr] (awgn) -- (gen)
    node[pos=0.5, above=2pt, ts] {$\hat{\mathbf{y}}$};
  \draw[arr] (gen) -- (out)
    node[pos=0.5, above=2pt, ts] {$1{\times}3{\times}256{\times}256$};

  % Discriminator 单独放下方，留足垂直间距
  \node[disc, fill=mPink!40, below=1.2cm of awgn] (disc) {Discriminator\\{\tiny 仅训练}};
  \draw[darr] (out.south) -- ++(0,-0.25) -| ++(-2.5,0) |- (disc.east);
  \draw[darr] (awgn.south) -- (disc.north);

  % === Encoder 微观（放在 Encoder 正下方，与 Disc 错开）===
  \node[micro, fill=mBlue!30, below=0.65cm of enc] (e1) {Conv};
  \node[micro, fill=mBlue!30, right=0.12cm of e1] (e2) {BN};
  \node[micro, fill=mBlue!30, right=0.12cm of e2] (e3) {ReLU};
  \node[micro, fill=mBlue!25, right=0.15cm of e3] (e4) {ResBlock$\times$2};
  \draw[arr, -] (e1) -- (e2) -- (e3) -- (e4);
  \draw[arr, dashed, mDark!60] (enc.south) -- (e2.north);

  % === Generator 微观（放在 Generator 正下方）===
  \node[micro, fill=mPurple!30, below=0.65cm of gen] (g1) {Conv};
  \node[micro, fill=mPurple!30, right=0.12cm of g1] (g2) {BN};
  \node[micro, fill=mPurple!30, right=0.12cm of g2] (g3) {ReLU};
  \node[micro, fill=mPurple!25, right=0.15cm of g3] (g4) {ResBlock$\times$5};
  \draw[arr, -] (g1) -- (g2) -- (g3) -- (g4);
  \draw[arr, dashed, mDark!60] (gen.south) -- (g2.north);

  % === 背景分组 ===
  \begin{scope}[on background layer]
    \node[fit=(img)(enc)(pwr), fill=mBlue!12, draw=mDark!25,
      rounded corners=4pt, inner sep=5pt, semithick,
      label={[font=\scriptsize\sffamily, color=mDark]above:发送端}] {};
    \node[fit=(gen)(out), fill=mPurple!12, draw=mDark!25,
      rounded corners=4pt, inner sep=5pt, semithick,
      label={[font=\scriptsize\sffamily, color=mDark]above:接收端}] {};
  \end{scope}

  % === 图例（放在右下，预留边距）===
  \node[legendbox, anchor=north east] at ([xshift=-3pt, yshift=-3pt]current bounding box.south east)
    [align=left, text width=3.8cm] {
    \textbf{图例}\\[1pt]
    \colorbox{mBlue!40}{\rule{0.15cm}{0.1cm}\hspace{0.06cm}} 编码 \quad
    \colorbox{mPurple!40}{\rule{0.15cm}{0.1cm}\hspace{0.06cm}} 解码 \quad
    \colorbox{mGreen!50}{\rule{0.15cm}{0.1cm}\hspace{0.06cm}} 信道\\
    实线：推理 \quad 虚线：训练
  };
\end{tikzpicture}
\end{document}
